\documentclass[UTF8, 11pt]{ctexart}

\usepackage[margin=1in]{geometry}
\usepackage{amsmath, amssymb}
\usepackage{booktabs}
\usepackage{hyperref}
\usepackage{url}
\usepackage{enumitem}
\setlist{leftmargin=*}
\setlength{\parskip}{0.35em}
\setlength{\parindent}{0pt}

\title{\textbf{阶段 2：代表性论文与方法选择}\\
用于平衡图划分的 Kernighan--Lin 细化方法}
\author{小组成员：姓名 1，姓名 2，姓名 3}
\date{CS240：算法设计与分析，2026 年春季学期}

\begin{document}

\maketitle

\section{算法基础铺垫}

阶段 1 已经把应用问题形式化为平衡图划分。在选择具体论文之前，需要先说明本
项目中几类方法背后的课内算法思想。

第一，平衡图划分是带约束的优化问题。目标是最小化割值，但划分还必须保持均衡。
这不同于随便找一个小割。一个割值很小的划分可能把几乎所有顶点都放到同一侧，
这样的结果无法用于多个 工作节点的任务分配。正是负载均衡约束，使这个问题更
接近算法课中讨论的困难图划分问题。

第二，本项目比较三类算法设计风格。贪心基线在构造划分时做局部看起来合理的
选择；谱二分把离散图问题松弛为图拉普拉斯矩阵上的线性代数问题；
Kernighan--Lin 使用局部搜索：从已有解出发，不断尝试通过交换顶点改进解。这些
方法不是互不相关的小技巧，而是处理困难离散优化问题时常见的三种思路。

第三，本阶段通过渐进复杂度分析让这些方法可以被比较。我们不只问哪个方法割值
更低，也要问它需要多少时间，以及代价如何依赖于 $n=|V|$、$m=|E|$ 和分区数
$k$。这为后续实现和实验阶段做铺垫。

\section{选择的代表性方法}

本项目选择复现的代表性方法是 Kernighan--Lin 图划分启发式算法。原论文研究的
主要问题是将图划分为两个大小相等的部分，同时最小化跨越割的边数或边权。该
方法是经典局部搜索算法：从一个初始划分出发，计算交换两个分区中顶点所带来的
增益，执行一系列可能有利的交换，并保留累计割值下降最大的交换前缀。

我们选择该方法，是因为它非常符合课程项目要求。它是经典算法过程，有清晰的
形式化问题定义，可以进行有意义的复杂度分析，并且能够自然连接到现代分布式
优化应用。在本项目中，Kernighan--Lin 方法被用作平衡 $k$ 路图划分的细化步骤。
由于原始 Kernighan--Lin 算法是二分算法，我们通过在当前不同分区对之间执行成对
细化，将其改写为 $k$ 路划分场景中的实用算法。

\section{正式问题定义}

设 $G=(V,E,w)$ 是一个无向加权图。对于任意二分 $(A,B)$，割值定义为
\[
    C(A,B) = \sum_{u\in A,\;v\in B,\;(u,v)\in E} w_{uv}.
\]
平衡二分问题要求求解
\[
    \min_{A,B} C(A,B)
    \quad
    \text{subject to}
    \quad
    A\cup B=V,\; A\cap B=\emptyset,\; |A|=|B|,
\]
当 $|V|$ 为奇数时，可以使用近似平衡版本。

本项目将问题推广到 $k$ 个分区。输出是映射
$p:V\to\{1,\ldots,k\}$，其中 $p(u)$ 表示顶点 $u$ 被分配到的 工作节点或分区。
$k$ 路割值定义为
\[
    \mathrm{cut}(p)
    =
    \sum_{(u,v)\in E,\; p(u)\neq p(v)} w_{uv}.
\]
负载均衡比例定义为
\[
    \rho(p)
    =
    \frac{\max_i |\{v\in V:p(v)=i\}|}{|V|/k}.
\]
当 $\rho(p)=1$ 时划分完全均衡；略大于 1 的值表示轻微不均衡。

\section{核心算法描述}

在图二分场景下，Kernighan--Lin 算法为每个顶点定义外部代价和内部代价。对于
顶点 $a\in A$，内部代价 $I(a)$ 是 $a$ 与同侧顶点之间边权之和，外部代价
$E(a)$ 是 $a$ 与另一侧顶点之间边权之和。差值
\[
    D(a) = E(a) - I(a)
\]
表示顶点 $a$ 被移动到另一侧的潜在收益。若交换 $a\in A$ 和 $b\in B$，交换增益为
\[
    g(a,b) = D(a) + D(b) - 2w_{ab},
\]
其中若 $a$ 与 $b$ 之间没有边，则 $w_{ab}=0$。正增益表示该交换可以降低割值。

算法的一轮 pass 通常如下：
\begin{enumerate}
    \item 从一个平衡划分 $(A,B)$ 出发；
    \item 将所有顶点标记为未锁定；
    \item 反复选择未锁定顶点对中增益最大的 $(a,b)$；
    \item 锁定被选择的顶点，并记录此次交换增益；
    \item 当所有可能顶点对都被选择后，寻找累计增益最大的交换前缀；
    \item 只有当最大累计增益为正时，才真正应用这一交换前缀。
\end{enumerate}

该方法的重要思想是：算法允许暂时选择局部看起来不好的交换，因为后续交换序列
可能带来整体正收益。因此，Kernighan--Lin 不是简单的贪心单步改进，而是带有
序列级别回看机制的局部搜索。

\section{到 $k$ 路划分的改写}

原始方法只处理两个集合。本项目在 $k$ 路框架中复用同样的局部搜索思想。首先，
我们用递归谱二分得到一个初始 $k$ 路划分。然后，对于每一对分区 $(i,j)$，取出
当前属于这两个分区的顶点诱导子图。接着，以当前两个分区作为初始二分，运行
成对 Kernighan--Lin 二分过程。若得到的局部更新不会增加全局割值，则接受该
更新。

这种改写适合课程项目，因为它保留了 Kernighan--Lin 的核心细化思想，同时让方法
可以应用到本项目中的 $k$ 个 工作节点分布式优化场景。它也便于与两个基线比较：
\begin{itemize}
    \item 贪心均衡划分基线，特点是简单且运行速度快；
    \item 递归谱二分，特点是使用线性代数松弛得到更强的初始划分。
\end{itemize}

\section{复杂度分析}

设 $n=|V|$，$m=|E|$，$k$ 是分区数。

\textbf{贪心基线。}
贪心方法逐个分配顶点。对每个顶点，它在可行分区中计算局部增量割值，并选择
代价较低的分区。直接实现的复杂度大约为 $O(mk)$，或者写作
$O(nk+\sum_v \deg(v)k)$，具体取决于是否维护增量代价。在本项目的小规模
benchmark 中，该方法很快，适合作为低运行时间基线。

\textbf{谱二分。}
递归谱二分在每次切分时构造图拉普拉斯矩阵并计算特征向量。若使用稠密矩阵特征
分解，规模为 $s$ 的子图一次切分复杂度为 $O(s^3)$。递归二分的总成本主要由较大
规模的切分决定。对于课程规模实验，这是可以接受的；若扩展到大图，则应使用
稀疏特征求解器。

\textbf{Kernighan--Lin 细化。}
对于含 $s$ 个顶点的二分子问题，朴素一轮 pass 可能需要评估大量顶点对，复杂度
约为 $O(s^2)$ 或更高，具体取决于数据结构。$k$ 路细化会在若干轮中考虑
$\binom{k}{2}$ 个分区对。若全局细化轮数为 $T$，每个成对子问题规模约为
$2n/k$，则粗略估计为
\[
    O\left(T \binom{k}{2}\left(\frac{2n}{k}\right)^2\right)
    =
    O(Tn^2),
\]
这里忽略了图稀疏性和具体库实现细节。实际运行时间会受到图密度、接受的细化
次数和停止条件影响。

\section{计划复现的实验目标}

本项目并不试图完整复现 Kernighan--Lin 原论文中的所有实验表格，而是复现该方法
最核心的算法行为：
\begin{enumerate}
    \item 从一个初始划分出发；
    \item 应用 Kernighan--Lin 风格的局部细化；
    \item 相比初始划分或简单基线降低割值；
    \item 保持分区大小基本均衡；
    \item 测量局部细化带来的额外运行时间。
\end{enumerate}

预期实验现象包括：
\begin{itemize}
    \item Kernighan--Lin 细化通常应比贪心划分得到更低割值；
    \item 它在许多实例上应能改进或持平谱划分，因为谱划分提供较强初始解，而
    局部搜索可以修补局部错误；
    \item 它的运行时间应高于贪心基线，因为需要额外细化 pass；
    \item 更低的割边权重应对应一致性仿真中更少的跨工作节点 消息。
\end{itemize}

\section{评估计划}

实验将在网格图、Erdos--Renyi 随机图、随机几何图和 Zachary karate club 图上
进行。这些实例分别覆盖结构化图、随机图、空间图和小型真实社交网络。评价指标
包括割边权重、归一化割、负载均衡比例、运行时间、一致性最终误差和跨分区消息
数量。

这一评估计划直接服务于项目目标：展示经典图划分启发式算法如何降低现代分布式
优化场景中的通信成本。

\begin{thebibliography}{9}

\bibitem{kernighan1970}
B. W. Kernighan and S. Lin.
\newblock An efficient heuristic procedure for partitioning graphs.
\newblock \emph{Bell System Technical Journal}, 1970.

\bibitem{karypis1998}
G. Karypis and V. Kumar.
\newblock A fast and high quality multilevel scheme for partitioning irregular graphs.
\newblock \emph{SIAM Journal on Scientific Computing}, 1998.

\end{thebibliography}

\end{document}

